\section{Sucesiones \texorpdfstring{$M$}{M}-regulares y sucesiones
\texorpdfstring{$\mathscr F$}{F}-regulares}
\label{section:0.15-es}

\subsection{Sucesiones \texorpdfstring{$M$}{M}-regulares y sucesiones
\texorpdfstring{$M$}{M}-cuasirregulares}
\label{subsection:0.15.1-es}
\label{0.15.1-es}

\begin{env}[15.1.1]
\label{0.15.1.1-es}
\oldpage[0\textsubscript{IV-1}]{12}
Sean $A$ un anillo y $N$ un $A$-módulo~; en lo que sigue, denotaremos
abreviadamente por $N[T_1,\ldots,T_n]$ el $A$-módulo graduado
$N\otimes_A A[T_1,\ldots,T_n]$ (siendo los $T_i$ indeterminadas), donde $N$
lleva la graduación trivial y $A[T_1,\ldots,T_n]$ la graduación por el grado total.

Sean $M$ un $A$-módulo y $f_i$ ($1\leq i\leq n$) elementos de $A$;
sea $\mathfrak{J}=\sum_{i=1}^n f_iA$ el ideal de $A$ engendrado por los $f_i$,
y dotemos a $M$ de la filtración $\mathfrak{J}$-preádica~; se define entonces un
\emph{homomorfismo graduado sobreyectivo de $A$-módulos graduados}, de grado $0$,
\[
  \varphi : (\gr_0(M))[T_1,\ldots,T_n]\longrightarrow\gr_\bullet(M)
  \tag{15.1.1.1}\label{0.15.1.1.1-es}
\]
de la manera siguiente~: si $\xi\in\gr_0(M)$ es la clase módulo
$\mathfrak{J}M$ de un elemento $x\in M$, al par formado por $\xi$ y un
polinomio $P[T_1,\ldots,T_n]$ homogéneo de grado $k$ se le hace corresponder la
clase módulo $\mathfrak{J}^{k+1}M$ de $P(f_1,\ldots,f_n)x$~; es inmediato
que esta clase
\oldpage[0\textsubscript{IV-1}]{13}
solo depende de $\xi$ y de $P$, y, si $S_k$ es el conjunto de los polinomios
homogéneos de grado $k$ de $A[T_1,\ldots,T_n]$, queda así definida una
aplicación $A$-lineal
\[
  \varphi_k:\gr_0(M)\otimes_A S_k\longrightarrow\gr_k(M),
\]
que es sobreyectiva por la definición de $\mathfrak{J}^kM$.

Diremos en lo que sigue que un elemento $f\in A$ es \emph{no divisor de
cero en un $A$-módulo $M$} si la homotecia $f_M:x\mapsto fx$ de $M$ es
\emph{inyectiva}.
\end{env}

\begin{lemma}[15.1.2]
\label{0.15.1.2-es}
Sean $M$ un $A$-módulo y $f$ un elemento de $A$, y dotemos a $M$ de la
filtración $(fA)$-preádica. Si $f$ es no divisor de cero en $M$,
el homomorfismo canónico definido en (15.1.1)~:
\[
  \varphi:(M/fM)[T]\longrightarrow\gr_\bullet(M)
  \tag{15.1.2.1}\label{0.15.1.2.1-es}
\]
es biyectivo. El recíproco es cierto si $M$ es separado para la topología
$(fA)$-preádica.
\end{lemma}

\begin{proof}
Se tiene aquí $\gr_k(M)=f^kM/f^{k+1}M$~; decir que $\varphi$ es inyectiva
significa que, para todo $x\in M$ y todo $k$, la relación
$f^kx\in f^{k+1}M$ implica $x\in fM$~; es claro que así ocurre si $f$ es no
divisor de cero en $M$. Recíprocamente, obsérvese que la homotecia
$f_M:x\mapsto fx$ define, por paso a los cocientes, homomorfismos
$\gr_k(M)\to\gr_{k+1}(M)$, y, por tanto, un homomorfismo graduado $g$ de grado
$1$ de $\gr_\bullet(M)$ en sí mismo~; decir que $\varphi$ es inyectivo expresa
también que $g$ es inyectivo. Si se supone separada la filtración de $M$, de
ello resulta que $f_M$ es inyectiva (Bourbaki, \emph{Alg.~comm.}, cap.~III,
\S~2, núm.~8, cor.~1 del teorema~1).
\end{proof}

\begin{corollary}[15.1.3]
\label{0.15.1.3-es}
Sean $A$ un anillo noetheriano y $M$ un $A$-módulo de tipo finito. Para que
$f$ sea no divisor de cero en $M$, es necesario y suficiente que $\varphi$ sea
biyectivo.
\end{corollary}

\begin{proof}
Solo hay que demostrar que la condición es suficiente~; para probar que $f_M$
es inyectiva, basta mostrar que, para todo ideal primo $\mathfrak p$ de $A$,
la homotecia de razón $f'=f/1\in A'=A_{\mathfrak p}$ en el
$A_{\mathfrak p}$-módulo $M_{\mathfrak p}=M'$ es inyectiva (Bourbaki,
\emph{Alg.~comm.}, cap.~II, \S~3, núm.~4, teorema~1). Ahora bien, se tiene
$f'^kM'/f'^{k+1}M'=(f^kM/f^{k+1}M)\otimes_A A'$, y es inmediato que
el homomorfismo canónico $\varphi':(M'/f'M')[T]\to\gr_\bullet(M')$ es igual a
$\varphi\otimes 1_{A'}$~; si $\varphi$ es inyectivo, lo mismo vale, por tanto,
para $\varphi'$ ($0_{\mathrm I}$, 1.3.2). Queda así reducido al caso en que $A$
es un anillo noetheriano \emph{local}, cuyo ideal maximal denotaremos por
$\mathfrak m$. Si $f\notin\mathfrak m$, entonces $f$ es invertible y no hay
nada que demostrar. Si $f\in\mathfrak m$, se sabe que $M$ es separado para la
topología $(fA)$-preádica ($0_{\mathrm I}$, 7.3.5), y basta, por tanto,
aplicar (15.1.2).
\end{proof}

\begin{definition}[15.1.4]
\label{0.15.1.4-es}
Sean $A$ un anillo y $M$ un $A$-módulo. Se dice que un elemento $f\in A$ es
\emph{$M$-regular} (resp. \emph{$M$-cuasirregular}) si $f$ es no divisor de
cero en $M$ (resp. si el homomorfismo (15.1.2.1) es biyectivo).
\end{definition}

Acabamos, pues, de demostrar que todo elemento $M$-regular es
$M$-cuasirregular, y que el recíproco es cierto si $A$ es noetheriano y $M$
un $A$-módulo de tipo finito. Cuando ya no se supone que $M$ sea de tipo
finito, este recíproco puede fallar~: por ejemplo, si $A$ es un anillo
principal que no es un cuerpo, $K$ su cuerpo de fracciones y $M$ el
$A$-módulo $K/A$, se tiene $fM=M$ para todo $f\neq0$ en $A$, por lo que los
dos miembros de (15.1.2.1) se reducen a $0$, pero $f$ es entonces divisor de
cero en $M$.

\begin{env}[15.1.5]
\label{0.15.1.5-es}
\oldpage[0\textsubscript{IV-1}]{14}
Supongamos ahora que $M$ está dotado de una \emph{filtración} $(M_k)$ tal que
$M_0=M$, y denotemos por
$\gr_\bullet(M)=\bigoplus_{k\geq0}M_k/M_{k+1}$ el $A$-módulo graduado asociado.
Con las notaciones de (15.1.1), pongamos, para todo $k\geq0$,
\[
  M'_k=M_k+\mathfrak{J}M_{k-1}+\cdots+
  \mathfrak{J}^{k-1}M_1+\mathfrak{J}^kM_0.
  \tag{15.1.5.1}\label{0.15.1.5.1-es}
\]
Es claro que $(M'_k)$ es una filtración sobre $M$. Cuando
$M_k=\mathfrak{Q}^kM$, donde $\mathfrak Q$ es un ideal de $A$, la filtración
$(M'_k)$ no es otra que la filtración
$(\mathfrak{J}+\mathfrak{Q})$-preádica. Denotaremos por
$\gr'_\bullet(M)$ el $A$-módulo graduado asociado a $(M'_k)$. Denotemos asimismo
\[
  (\gr_\bullet(M)\otimes_A(A/\mathfrak{J}))[T_1,\ldots,T_n]
\]
el producto tensorial de $A$-módulos graduados ($\mathrm{II}$, 2.1.2)
\[
  (\gr_\bullet(M)\otimes_A(A/\mathfrak{J}))
  \otimes_A A[T_1,\ldots,T_n].
\]
Se define un \emph{homomorfismo graduado sobreyectivo de $A$-módulos graduados},
de grado $0$,
\[
  \psi:(\gr_\bullet(M)\otimes_A(A/\mathfrak{J}))[T_1,\ldots,T_n]
  \longrightarrow\gr'_\bullet(M)
  \tag{15.1.5.2}\label{0.15.1.5.2-es}
\]
de la manera siguiente~: si $\alpha\in A/\mathfrak{J}$ es la clase de un
elemento $a\in A$, $\xi\in M_k/M_{k+1}$ la clase de un elemento $x\in M_k$ y
$P[T_1,\ldots,T_n]$ un polinomio homogéneo de grado $h$, se hace corresponder
a la terna $(\alpha,\xi,P)$ la clase módulo $M'_{h+k+1}$ de
$P(f_1,\ldots,f_n)ax$~; se comprueba de nuevo inmediatamente que este elemento
solo depende de $\xi$, $\alpha$ y $P$, y esto define el homomorfismo graduado
(15.1.5.2), que es sobreyectivo en virtud de la definición (15.1.5.1). Se
recupera el homomorfismo (15.1.1.1) considerando la filtración sobre $M$ tal
que $M_1=0$.
\end{env}

\begin{lemma}[15.1.6]
\label{0.15.1.6-es}
Se conservan las notaciones de (15.1.5) y se supone $n=1$, poniendo $f=f_1$.
Si $f$ es $\gr_\bullet(M)$-regular, el homomorfismo (15.1.5.2) es biyectivo.
\end{lemma}

\begin{proof}
Sea $Q_k$ (resp. $Q'_k$) el sub-$A$-módulo de los términos de grado $k$ en
el primer (resp. segundo) miembro de (15.1.5.2). Dotemos a $Q_k$ de la filtración
\[
  (Q_k)_i=\sum_{j\leq k-i}
  (\gr_{k-j}(M)\otimes_A(A/fA))T^j
  \tag{15.1.6.1}\label{0.15.1.6.1-es}
\]
de modo que $(Q_k)_0=Q_k$ y $(Q_k)_{k+1}=0$, y, para esta filtración, se tiene
\[
  \gr_i(Q_k)=(\gr_i(M)\otimes_A(A/fA))T^{k-i}.
\]
Dotemos a $Q'_k$ de la filtración formada por los
$\psi((Q_k)_i)=(Q'_k)_i$. Como estas filtraciones son \emph{finitas}, basta
probar que los homomorfismos $\psi_{ki}:\gr_i(Q_k)\to\gr_i(Q'_k)$ deducidos de
$\psi_k$ son inyectivos (Bourbaki, \emph{Alg.~comm.}, cap.~III, \S~3,
núm.~8, cor.~1 del teorema~1). Ahora bien, se tiene
\[
  \gr_i(Q_k)=((M_i/M_{i+1})\otimes_A(A/fA))T^{k-i}~;
\]
por otro lado, $(Q'_k)_{i+1}$ es la imagen de
\[
  M_k+fM_{k-1}+\cdots+f^{k-i-1}M_{i+1}
\]
\oldpage[0\textsubscript{IV-1}]{15}
en $M'_k/M'_{k+1}$. Nos reducimos, por tanto, a demostrar que, para $x\in M_i$,
la relación
\[
  f^{k-i}x\in M_k+fM_{k-1}+\cdots+
  f^{k-i-1}M_{i+1}+M'_{k+1}
  \tag{15.1.6.2}\label{0.15.1.6.2-es}
\]
implica $x\in fM_i+M_{i+1}$. Ahora bien, el segundo miembro de (15.1.6.2) está
contenido en $M_{i+1}+M'_{k+1}$, y, en virtud de (15.1.5.1), $M'_{k+1}$ está
contenido en $M_{i+1}+f^{k+1-i}M$. La hipótesis sobre $f$ implica que $f$ es
no divisor de cero en $M/M_h$ para todo $h$ (Bourbaki, \emph{loc.~cit.}). La
relación $f^{k-i}x\in f^{k+1-i}M+M_{i+1}$ implica, pues (razonando en
$M/M_{i+1}$), la existencia de un $y\in M$ tal que $x-fy\in M_{i+1}$, y como
$x\in M_i$, se tiene $fy\in M_i$, de donde (razonando en $M/M_i$) se sigue
$y\in M_i$, y finalmente $x\in fM_i+M_{i+1}$. C.Q.D.
\end{proof}

\begin{definition}[15.1.7]
\label{0.15.1.7-es}
Sea $(f_i)_{1\leq i\leq n}$ una sucesión de elementos de $A$. Se dice que es
\emph{$M$-regular} si, para $1\leq i\leq n$, $f_i$ es
$\left(M/(\sum_{1\leq j\leq i-1}f_jM)\right)$-regular. Se dice que la sucesión
$(f_i)$ es \emph{$M$-cuasirregular} si el homomorfismo canónico $\varphi$
definido en (15.1.1.1) es biyectivo.

Cuando $M=A$, se dice simplemente «~sucesión regular~» y
«~sucesión cuasirregular~».
\end{definition}

\begin{proposition}[15.1.8]
\label{0.15.1.8-es}
Con las notaciones de (15.1.5), se supone que la sucesión
$(f_i)_{1\leq i\leq n}$ es $\gr_\bullet(M)$-regular. Entonces
el homomorfismo canónico (15.1.5.2) es biyectivo.
\end{proposition}

\begin{proof}
La proposición no es sino (15.1.6) para $n=1$~; razonemos por inducción sobre
$n$. Pongamos
\[
  \mathfrak{J}''=\sum_{i=1}^{n-1}f_iA,
\]
de modo que $\mathfrak{J}=\mathfrak{J}''+f_nA$, y sea
\[
  M''_k=M_k+\mathfrak{J}''M_{k-1}+\cdots+
  \mathfrak{J}''^{k-1}M_1+\mathfrak{J}''^kM_0~;
\]
se puede escribir, por tanto,
\[
  M'_k=M''_k+f_nM''_{k-1}+\cdots+f_n^{k-1}M''_1+f_n^kM''_0.
\]
Denotemos por $\gr''_\bullet(M)$ el $A$-módulo graduado asociado a la filtración
$(M''_k)$ de $M$. Salvo isomorfismos canónicos, se puede escribir
\[
  (\gr_\bullet(M)\otimes(A/\mathfrak{J}))\otimes_A A[T_1,\ldots,T_n]
  =
  (\gr_\bullet(M)\otimes(A/\mathfrak{J}''))[T_1,\ldots,T_{n-1}]
  \otimes(A/f_nA)\otimes_A A[T_n]
\]
y, por consiguiente, (15.1.5.2) se factoriza en
\[
  (\gr''_\bullet(M)\otimes(A/f_nA))[T_n]\longrightarrow\gr'_\bullet(M)
  \tag{15.1.8.1}\label{0.15.1.8.1-es}
\]
y
\[
  ((\gr_\bullet(M)\otimes(A/\mathfrak{J}''))[T_1,\ldots,T_{n-1}])
  \otimes(A/f_nA)\otimes_A A[T_n]
  \xrightarrow{\psi'\otimes1}
  \gr''_\bullet(M)\otimes(A/f_nA)\otimes_A A[T_n]
\]
donde $\psi'$ es la aplicación (15.1.5.2) en la que se reemplazan $n$ por
$n-1$, $\mathfrak{J}$ por $\mathfrak{J}''$ y $\gr'_\bullet(M)$ por
$\gr''_\bullet(M)$. La hipótesis de inducción implica que $\psi'$ es
biyectiva, y queda, por tanto, comprobar que lo mismo vale para (15.1.8.1).
Aplicando (15.1.6), basta mostrar que $f_n$ es
$\gr''_\bullet(M)$-regular. Pero $\gr''_\bullet(M)$ se identifica con
\[
  (\gr_\bullet(M)\otimes(A/\mathfrak{J}''))
  [T_1,\ldots,T_{n-1}]~;
\]
por hipótesis, la homotecia de razón $f_n$ en
$\gr_\bullet(M)\otimes(A/\mathfrak{J}'')$ es inyectiva~; por consiguiente, lo
mismo vale para la homotecia de razón $f_n$ en $\gr''_\bullet(M)$, puesto que
$A[T_1,\ldots,T_{n-1}]$ es un $A$-módulo libre.
\end{proof}
\begin{proposition}[15.1.9]
\label{0.15.1.9-es}
Sean $M$ un $A$-módulo y $(f_i)_{1\leq i\leq n}$ una sucesión de elementos de
$A$. Si la sucesión $(f_i)$ es $M$-regular, es $M$-cuasirregular. El recíproco
es cierto cuando se supone además que los $A$-módulos
\[
  M,\quad M/f_1M,\quad\ldots,\quad
  M/\left(\sum_{1\leq j\leq n-1}f_jM\right)
\]
son separados para la topología $\mathfrak{J}$-preádica (donde
$\mathfrak{J}=\sum_{i=1}^n f_iA$).
\end{proposition}

\begin{proof}
La primera afirmación resulta de (15.1.8), tomando $M_1=0$, y, por
consiguiente, el homomorfismo (15.1.5.2) se reduce a (15.1.1.1).
Recíprocamente, supongamos que el homomorfismo $\varphi$ de (15.1.1.1) es
biyectivo. Obsérvese que, si se pone
\[
  \mathfrak{J}''=\sum_{i=1}^{n-1}f_iA,
\]
\oldpage[0\textsubscript{IV-1}]{16}
y si se denota por $\gr''_\bullet(M)$ el $A$-módulo graduado asociado a $M$
filtrado por la filtración $\mathfrak{J}''$-preádica, $\varphi$ se factoriza en
\[
  (\gr''_\bullet(M)\otimes(A/f_nA))[T_n]\longrightarrow\gr_\bullet(M)
  \tag{15.1.9.1}\label{0.15.1.9.1-es}
\]
y en $\varphi'\otimes1$, donde $\varphi'$ es el homomorfismo canónico
\[
  (\gr''_0(M))[T_1,\ldots,T_{n-1}]
  \longrightarrow\gr''_\bullet(M).
\]
Como estos dos homomorfismos son sobreyectivos, son biyectivos si se supone que
$\varphi$ es biyectivo; por tanto, $\varphi'$ es necesariamente inyectivo y,
por consiguiente, biyectivo, puesto que es sobreyectivo. Se puede entonces
razonar por recurrencia sobre $n$ (el caso $n=1$ resulta de (15.1.2)), pues
sobre
\[
  M,\quad M/f_1M,\quad\ldots,\quad
  M/\left(\sum_{j=1}^{n-2}f_jM\right)
\]
la topología $\mathfrak{J}''$-preádica es más fina que la topología
$\mathfrak{J}$-preádica y, por consiguiente, es separada. Se ve, pues, que la
sucesión $(f_i)_{1\leq i\leq n-1}$ es $M$-regular. Por otra parte, demostremos
que la hipótesis implica que, si $y\in M/\mathfrak{J}''M$ es tal que
\[
  f_n^ky\in f_n^{k+1}(M/\mathfrak{J}''M),
\]
entonces $y\in f_n(M/\mathfrak{J}''M)$. En efecto, si $x$ es un representante
de $y$, demostremos que $f_n^kx\in\mathfrak{J}^{k+1}M$; nuestra afirmación
resultará entonces del hecho de que (15.1.9.1) es inyectivo. La hipótesis
implica
\[
  f_n^kx\in f_n^{k+1}M+\mathfrak{J}''M\subset\mathfrak{J}M,
\]
de donde, como $\varphi$ es inyectivo,
$f_n^{k-1}x\in\mathfrak{J}^kM$; por recurrencia se obtiene así
$x\in\mathfrak{J}M$, y de ahí, finalmente,
$f_n^kx\in\mathfrak{J}^{k+1}M$. Se termina el razonamiento como en
(15.1.2): sobre $M/\mathfrak{J}''M$, la filtración
$\mathfrak{J}$-preádica es idéntica a la filtración
$(f_nA)$-preádica, luego es separada por hipótesis, y se concluye que la
homotecia de razón $f_n$ en $M/\mathfrak{J}''M$ es inyectiva, lo que termina
de demostrar que la sucesión $(f_i)_{1\leq i\leq n}$ es $M$-regular.
\end{proof}

\begin{corollary}[15.1.10]
\label{0.15.1.10-es}
Con las notaciones de (15.1.9), supongamos que, para todo $A$-módulo $N$ de
tipo finito, $M\otimes_A N$ es separado para la topología
$\mathfrak{J}$-preádica. Entonces, para que la sucesión $(f_i)$ sea
$M$-regular, es necesario y suficiente que sea $M$-cuasirregular.
\end{corollary}

Obsérvese que la hipótesis de (15.1.10) implica que, si la sucesión
$(f_i)_{1\leq i\leq n}$ es $M$-regular, también lo es la sucesión
$(f_{\sigma(i)})$ para toda permutación $\sigma$ del intervalo $[1,n]$.

\begin{corollary}[15.1.11]
\label{0.15.1.11-es}
Sean $A$ un anillo noetheriano, $M$ un $A$-módulo de tipo finito y
$(f_i)_{1\leq i\leq n}$ una sucesión finita de elementos contenidos en el
radical de $A$. Entonces, para que la sucesión $(f_i)$ sea $M$-regular, es
necesario y suficiente que sea $M$-cuasirregular.
\end{corollary}

\begin{proof}
Es un caso particular de (15.1.10), puesto que todo $A$-módulo de tipo finito
es entonces separado para la topología $\mathfrak{J}$-preádica
($0_{\mathrm I}$, 7.3.5).
\end{proof}

\begin{remarks}[15.1.12]
\label{0.15.1.12-es}
\mbox{}
\begin{enumerate}
  \item[(i)] Cuando $A$ es un anillo local noetheriano, $M$ un $A$-módulo de
  tipo finito, y cuando los $f_i$ están contenidos en el ideal maximal de $A$,
  la noción de sucesión $M$-regular fue introducida por J.-P.~Serre con el
  nombre de $M$-\emph{sucesión} [17].
  \item[(ii)] Incluso si $A$ es noetheriano y $M=A$, una sucesión $(f,g)$ de
  dos elementos de $A$ que es cuasirregular no es necesariamente regular. Por
  ejemplo, la sucesión $(0,1)$ es cuasirregular, puesto que entonces
  $\mathfrak{J}=fA+gA=A$, luego $\gr_\bullet(A)=0$; pero, si $A\neq0$, no es
  regular; obsérvese, no obstante, que en este caso la sucesión $(1,0)$ es
  regular, pues $1$ no es divisor de $0$ en $A$, y como $A/A=0$, $0$ no es
  divisor de cero en este módulo (cf. (15.1.1)).
  \item[(iii)]
  \oldpage[0\textsubscript{IV-1}]{17}
  En (15.1.6), en general no se puede sustituir la hipótesis de que $f$ es
  $\gr_\bullet(M)$-regular por la hipótesis de que es
  $\gr_\bullet(M)$-cuasirregular. Para verlo, retomemos el ejemplo (15.1.4),
  donde $A$ es un anillo principal de cuerpo de fracciones $K\neq A$, y
  tomemos $M=K$, definiéndose la filtración $(M_k)$ de $M$ por
  $M_0=M=K$, $M_1=A$, $M_k=0$ para $k\geq2$, de donde
  $\gr_0(M)=K/A$, $\gr_1(M)=A$, $\gr_k(M)=0$ para $k\geq2$. Si $f\neq0$ en
  $A$, se ha visto que $f$ no es $\gr_\bullet(M)$-regular, pero, como se tiene
  \[
    f^k\gr_\bullet(M)/f^{k+1}\gr_\bullet(M)=f^kA/f^{k+1}A,
  \]
  es inmediato que el homomorfismo (15.1.2.1) (donde $M$ se sustituye por
  $\gr_\bullet(M)$) es inyectivo; por tanto, $f$ es
  $\gr_\bullet(M)$-cuasirregular. Pero, con las notaciones de (15.1.6), se
  tiene entonces $M'_k=K$ para todo $k\geq0$, luego
  $\gr'_\bullet(M)=0$, mientras que el miembro izquierdo de (15.1.5.2) no se
  reduce a $0$.
\end{enumerate}
\end{remarks}

\begin{proposition}[15.1.13]
\label{0.15.1.13-es}
Sean $A$ un anillo, $M$ y $N$ dos $A$-módulos, y
$(f_i)_{1\leq i\leq n}$ una sucesión de elementos de $A$. Si la sucesión
$(f_i)$ es $M$-regular y si $N$ es un $A$-módulo plano, entonces la sucesión
$(f_i)$ es $M\otimes_A N$-regular. El recíproco es cierto si $N$ es un
$A$-módulo fielmente plano.
\end{proposition}

\begin{proof}
En efecto (sin hipótesis sobre $N$),
\[
  (M\otimes_A N)/
  \left(\sum_{j=1}^{i-1}f_j(M\otimes_A N)\right)
\]
es isomorfo a
\[
  \left(M/\left(\sum_{j=1}^{i-1}f_jM\right)\right)\otimes_A N~;
\]
basta entonces aplicar ($0_{\mathrm I}$, 6.1.1) y
($0_{\mathrm I}$, 6.4.1) a la homotecia de razón $f_i$ en
$M/(\sum_{j=1}^{i-1}f_jM)$.
\end{proof}

\begin{corollary}[15.1.14]
\label{0.15.1.14-es}
Sean $A$ un anillo, $M$ un $A$-módulo y $(f_i)_{1\leq i\leq n}$ una sucesión
finita de elementos de $A$. Sean $\rho:A\to B$ un homomorfismo de anillos,
$M'=M_{(B)}=M\otimes_A B$ y $f'_i=\rho(f_i)$ para $1\leq i\leq n$. Si $B$ es
un $A$-módulo plano y si la sucesión $(f_i)$ es $M$-regular, entonces la
sucesión $(f'_i)$ es $M'$-regular; el recíproco es cierto si $B$ es un
$A$-módulo fielmente plano.
\end{corollary}

\begin{proof}
En efecto,
\[
  M'/\left(\sum_{j=1}^{i-1}f'_jM'\right)
\]
se identifica con
\[
  \left(M/\left(\sum_{j=1}^{i-1}f_jM\right)\right)\otimes_A B,
\]
y la homotecia de razón $f'_i$ en el $B$-módulo
$M'/(\sum_{j=1}^{i-1}f'_jM')$ se identifica con la homotecia de razón $f_i$
en el $A$-módulo
\[
  \left(M/\left(\sum_{j=1}^{i-1}f_jM\right)\right)\otimes_A B~;
\]
la afirmación resulta, por tanto, de (15.1.13).
\end{proof}

\begin{proposition}[15.1.15]
\label{0.15.1.15-es}
Sean $A$ un anillo, $B$ una $A$-álgebra conmutativa,
$(f_i)_{1\leq i\leq n}$ una sucesión finita de elementos de $B$ y $M$ un
$B$-módulo. Supongamos que la sucesión $(f_i)$ es $M$-regular y que cada uno
de los $A$-módulos
\[
  M/\left(\sum_{j=1}^{i-1}f_jM\right)\quad(1\leq i\leq n)
\]
es plano. Entonces, para todo homomorfismo de anillos $\rho:A\to A'$, si se
pone $B'=B\otimes_A A'$, $M'=M\otimes_A A'$ y $f'_i=f_i\otimes1$
($1\leq i\leq n$), la sucesión $(f'_i)$ es $M'$-regular y los
$A'$-módulos
\[
  M'/\left(\sum_{j=1}^{i-1}f'_jM'\right)\quad(1\leq i\leq n)
\]
son planos.
\end{proposition}

\begin{proof}
Consideremos la sucesión (exacta por hipótesis)
\[
  0\longrightarrow M\xrightarrow{f_1}M\longrightarrow M/f_1M
  \longrightarrow0
\]
(siendo la flecha $M\xrightarrow{f_1}M$ la homotecia de razón $f_1$ en
$M$). La hipótesis de que $M/f_1M$ es $A$-plano implica que la sucesión
\[
  0\longrightarrow M\otimes_A A'
  \xrightarrow{f_1\otimes1}M\otimes_A A'
  \longrightarrow(M/f_1M)\otimes_A A'\longrightarrow0
\]
\oldpage[0\textsubscript{IV-1}]{18}
es exacta ($0_{\mathrm I}$, 6.1.2). Esto muestra que la homotecia de razón
$f'_1$ en $M'$ es inyectiva.

Pongamos a continuación, para $1\leq i\leq n-1$,
\[
  M_i=M/\left(\sum_{j=1}^i f_jM\right),\qquad
  M'_i=M'/\left(\sum_{j=1}^i f'_jM'\right)~;
\]
se tiene $M'_i=M_i\otimes_A A'$, $M_i/f_{i+1}M_i=M_{i+1}$,
$M'_i/f'_{i+1}M'_i=M'_{i+1}$; basta sustituir en el razonamiento anterior
$M$ y $f_1$ por $M_i$ y $f_{i+1}$ para concluir que la homotecia de razón
$f'_{i+1}$ en $M'_i$ es inyectiva, en virtud de la planitud de
$M_i/f_{i+1}M_i$. La última afirmación resulta de ($0_{\mathrm I}$, 6.2.1).
\end{proof}
\begin{proposition}[15.1.16]
\label{0.15.1.16-es}
Sean $A$, $B$ dos anillos locales noetherianos, $k$ el cuerpo residual de
$A$, $\varphi:A\to B$ un homomorfismo local y $M$ un $B$-módulo de tipo finito.
Sea $(f_i)_{1\leq i\leq n}$ una sucesión de elementos del ideal máximo de $B$
y, para cada $i$, sea $g_i$ la imagen de $f_i$ en $B\otimes_A k$. Las
condiciones siguientes son equivalentes:
\begin{enumerate}
  \item[\rm(a)] La sucesión $(f_i)$ es $M$-regular y los cocientes
  \[
    M_i=M/\left(\sum_{j=1}^i f_jM\right)
  \]
  son $A$-módulos planos para $1\leq i\leq n$.
  \item[\rm(b)] La sucesión $(f_i)$ es $M$-regular y el $A$-módulo $M_n$ es
  plano.
  \item[\rm(c)] El $A$-módulo $M$ es plano y la sucesión
  $(g_i)_{1\leq i\leq n}$ es $(M\otimes_A k)$-regular.
  \item[\rm(d)] El $A$-módulo $M$ es plano y, para todo homomorfismo de
  anillos $\rho:A\to A'$, si se pone $B'=B\otimes_A A'$,
  $M'=M\otimes_A A'$, $f'_i=f_i\otimes1$ ($1\leq i\leq n$), la sucesión
  $(f'_i)$ es $M'$-regular.
\end{enumerate}
\end{proposition}

\begin{proof}
Es claro que a) implica b); para demostrar que b) implica d), basta, en virtud
de (15.1.15), demostrar que b) implica que los $M_i$ son $A$-módulos planos
para $0\leq i\leq n-1$ (poniendo $M_0=M$). Como
$M_{i+1}=M_i/f_{i+1}M_i$, por inducción descendente nos reducimos a probar que,
si $h$ es un elemento del ideal máximo de $B$ tal que la multiplicación por
$h$ en $M$ sea inyectiva y que $M/hM$ sea un $A$-módulo plano, entonces $M$
es un $A$-módulo plano; pero esto no es sino
($0_{\mathrm{III}}$, 10.2.7). Es claro que c) es el caso particular de d)
obtenido tomando $A'=k$. Queda, pues, probar que c) implica a). Como $M$ es un
$A$-módulo plano y un $B$-módulo de tipo finito, se deduce de c) y de
($0_{\mathrm{III}}$, 10.2.4) que la multiplicación por $f_1$ en $M$ es
inyectiva y que su conúcleo $M/f_1M=M_1$ es un $A$-módulo plano. Razonemos por
inducción sobre $i$ y supongamos demostrado que $M_i$ es un $A$-módulo plano;
como la multiplicación por $g_{i+1}$ en $M_i\otimes_A k$ es inyectiva por
hipótesis, se deduce de nuevo de ($0_{\mathrm{III}}$, 10.2.4) que la
multiplicación por $f_{i+1}$ en $M_i$ es inyectiva y que
$M_i/f_{i+1}M_i=M_{i+1}$ es un $A$-módulo plano, lo que termina la
demostración. Se observará que la demostración de que c) implica a) no utiliza
el hecho de que los $f_i$ pertenezcan al ideal máximo de $B$.
\end{proof}

\begin{corollary}[15.1.17]
\label{0.15.1.17-es}
Bajo las hipótesis sobre $A$ y $B$ de (15.1.16), sean $M$ un $B$-módulo de
tipo finito, plano sobre $A$, y $R$ un sub-$B$-módulo de $M$ tal que $N=M/R$
sea un $A$-módulo plano. Sea $(f_i)_{1\leq i\leq n}$ una sucesión de elementos
del ideal máximo de $B$ que tenga las propiedades siguientes:
\begin{enumerate}
  \item[(i)] $f_iM\subset R$ para $1\leq i\leq n$.
  \item[(ii)] Si $g_i$ es la imagen de $f_i$ en $B\otimes_A k$, la sucesión
  $(g_i)_{1\leq i\leq n}$ es $(M\otimes_A k)$-regular.
  \item[(iii)] La suma de los sub-$(B\otimes_A k)$-módulos
  $g_i(M\otimes_A k)$ de $M\otimes_A k$ es la imagen canónica de
  $R\otimes_A k$ en $M\otimes_A k$ (igual al núcleo del homomorfismo canónico
  $M\otimes_A k\to N\otimes_A k$).
\end{enumerate}
Entonces la sucesión $(f_i)$ es $M$-regular y se tiene
$R=\sum_{i=1}^n f_iM$.
\end{corollary}

\begin{proof}
\oldpage[0\textsubscript{IV-1}]{19}
Como $N$ es un $A$-módulo plano, se deduce de la exactitud de la sucesión
$0\to R\to M\to N\to0$ que la sucesión
\[
  0\longrightarrow R\otimes_A k\longrightarrow M\otimes_A k
  \longrightarrow N\otimes_A k\longrightarrow0
\]
es exacta ($0_{\mathrm I}$, 6.1.2); se tiene, pues,
$R\otimes_A k=\sum_i g_i(M\otimes_A k)$. Como $M$ es un $B$-módulo de tipo
finito y $B$ es noetheriano, $R$ es también un $B$-módulo de tipo finito, de
donde $R=\sum_i f_iM$ en virtud del lema de Nakayama, pues los $f_i$ pertenecen
al ideal máximo de $B$.
\end{proof}

\begin{lemma}[15.1.18]
\label{0.15.1.18-es}
Sean $A$ un anillo, $M$ un $A$-módulo provisto de una filtración separada y
exhaustiva $(M_k)_{k\geq0}$, y sea $f\in A$ un elemento
$\gr_\bullet(M)$-regular. Entonces:
\begin{enumerate}
  \item[(i)] $f$ es $M$-regular.
  \item[(ii)] Para todo $k$, se tiene $M_k\cap fM=fM_k$, y la imagen canónica
  $M'_k$ de $M_k$ en $M'=M/fM$ es, por tanto, canónicamente isomorfa a
  $M_k/fM_k$.
  \item[(iii)] La sucesión
  \[
    0\longrightarrow M_{k+1}/fM_{k+1}\longrightarrow M_k/fM_k
    \longrightarrow\gr_k(M)/f\cdot\gr_k(M)\longrightarrow0
  \]
  es exacta y, si se dota a $M'$ de la filtración cociente $(M'_k)$,
  $\gr_\bullet(M')$ es, por tanto, canónicamente isomorfo a
  $\gr_\bullet(M)/f\cdot\gr_\bullet(M)$.
\end{enumerate}
\end{lemma}

\begin{proof}
La afirmación (i) se deduce de la hipótesis sobre $f$ y de Bourbaki,
\emph{Alg.~comm.}, cap.~III, \S~2, n.º~8, cor.~1 del teo.~1. Como
$\gr_\bullet(M/M_k)$ es un sumando directo de $\gr_\bullet(M)$, $f$ es
$\gr_\bullet(M/M_k)$-regular para todo $k$, luego $(M/M_k)$-regular por (i),
lo que establece (ii). Para probar (iii), basta ver que la aplicación
$M_{k+1}/fM_{k+1}\to M_k/fM_k$ es inyectiva: ahora bien, si
$x\in M_{k+1}$ es tal que $x\in fM_k$, se deduce de (ii) que
$x\in fM_{k+1}$, lo que termina de probar el lema.
\end{proof}

\begin{proposition}[15.1.19]
\label{0.15.1.19-es}
Sean $A$ un anillo, $M$ un $A$-módulo provisto de una filtración
$(M_k)_{k\geq0}$ tal que $M_0=M$, y $(f_i)_{1\leq i\leq n}$ una sucesión
$\gr_\bullet(M)$-regular de elementos de $A$. Supongamos además que, para
$0\leq i\leq n-1$, la filtración cociente de $(M_k)$ sobre
\[
  M/\left(\sum_{j\leq i}f_jM\right)
\]
es separada. Entonces la sucesión $(f_i)$ es $M$-regular y, si se pone
\[
  M'=M/\left(\sum_{i\leq n}f_iM\right)
\]
y se dota a $M'$ de la filtración cociente de $(M_k)$,
$\gr_\bullet(M')$ es canónicamente isomorfo a
\[
  \gr_\bullet(M)/
  \left(\sum_{i\leq n}f_i\gr_\bullet(M)\right).
\]
\end{proposition}

\begin{proof}
Para $n=1$, la proposición se deduce del lema (15.1.18). Razonemos por
inducción sobre $n$, poniendo
\[
  N=M/\left(\sum_{i\leq n-1}f_iM\right)
\]
y denotando por $(N_k)$ la filtración cociente de $(M_k)$ sobre $N$, que es
separada por hipótesis; además, $\gr_\bullet(N)$ es isomorfo a
\[
  \gr_\bullet(M)/
  \left(\sum_{i\leq n-1}f_i\gr_\bullet(M)\right).
\]
Por hipótesis, $f_n$ es, pues, $\gr_\bullet(N)$-regular, y se deduce del lema
(15.1.18), aplicado a $N$, que $f_n$ es $N$-regular y $\gr_\bullet(M')$ es
isomorfo a $\gr_\bullet(N)/f_n\gr_\bullet(N)$, lo que demuestra la
proposición.
\end{proof}

La proposición (15.1.19) se aplicará, en particular, a las filtraciones que
verifiquen una u otra de las hipótesis siguientes:
\begin{enumerate}
  \item[$1^\circ$] La filtración $(M_k)$ es finita y separada (pues esto
  implica que $M_k=0$ para $k$ suficientemente grande);
  \item[$2^\circ$] $A$ es un anillo noetheriano, $\mathfrak{J}$ un ideal
  contenido en el radical de $A$, $M$ un $A$-módulo de tipo finito y $(M_k)$
  la filtración $\mathfrak{J}$-preádica ($0_{\mathrm I}$, 7.3.5).
\end{enumerate}

\begin{corollary}[15.1.20]
\label{0.15.1.20-es}
\oldpage[0\textsubscript{IV-1}]{20}
Sean $A$ un anillo, $M$ un $A$-módulo y $(f_i)_{1\leq i\leq n}$ una sucesión
$M$-regular de elementos de $A$. Entonces, cualesquiera que sean los enteros
$\nu_i\geq1$, la sucesión $(f_i^{\nu_i})$ es $M$-regular y
$M/(\sum_{i\leq n}f_i^{\nu_i}M)$ posee una serie de composición cuyos
cocientes son isomorfos a $M/(\sum_{i\leq n}f_iM)$.
\end{corollary}

\begin{proof}
Razonemos por inducción sobre $n$; para $n=1$, las afirmaciones son inmediatas,
pues $f^jM/f^{j+1}M$ es isomorfo a $M/fM$, ya que la multiplicación por $f$ en
$M$ es inyectiva por hipótesis. Pongamos entonces
\[
  N=M/\left(\sum_{i\leq n-1}f_i^{\nu_i}M\right)~;
\]
por la hipótesis de inducción, $N$ admite una filtración finita $(N_k)$ para la
cual los $\gr_k(N)$ son isomorfos a
$M/(\sum_{i\leq n-1}f_iM)$; por hipótesis, $f_n$ es, pues,
$\gr_\bullet(N)$-regular; se deduce entonces de (15.1.19) que $f_n$ es
$N$-regular. Aplicando el caso $n=1$ demostrado al principio, se concluye que
$f_n^{\nu_n}$ es $N$-regular y que $N/f_n^{\nu_n}N$ admite una serie de
composición cuyos cocientes son isomorfos a $N/f_nN$; pero, para la filtración
de $N/f_nN$ cociente de la filtración $(N_k)$,
$\gr_k(N/f_nN)$ es isomorfo a $\gr_k(N)/f_n\gr_k(N)$ por (15.1.19) para todo
$k$, luego a $M/(\sum_{i\leq n}f_iM)$, lo que prueba el corolario.
\end{proof}

\begin{proposition}[15.1.21]
\label{0.15.1.21-es}
Sean $A$, $B$ dos anillos locales noetherianos, $\varphi:A\to B$ un
homomorfismo local, $M$ un $B$-módulo de tipo finito y
$(f_i)_{1\leq i\leq n}$ una sucesión $A$-regular formada por elementos del
ideal máximo de $A$. Las propiedades siguientes son equivalentes:
\begin{enumerate}
  \item[a)] $M$ es un $A$-módulo plano.
  \item[b)] Si $\mathfrak{J}$ es el ideal $\sum_{i\leq n}f_iA$,
  $M/\mathfrak{J}M$ es un $A/\mathfrak{J}$-módulo plano y la sucesión $(f_i)$
  es $M$-regular.
\end{enumerate}
\end{proposition}

\begin{proof}
El hecho de que a) implica b) se deduce de (15.1.13) y de
($0_{\mathrm{III}}$, 10.2.2). Recíprocamente, si la sucesión $(f_i)$ es
$M$-regular, es $M$-cuasi-regular (15.1.9), luego el homomorfismo
(15.1.1.1) es biyectivo. Pero, como también se supone que la sucesión
$(f_i)_{1\leq i\leq n}$ es $A$-regular, es $A$-cuasi-regular (15.1.9), y el
homomorfismo canónico (15.1.1.1) correspondiente
$(A/\mathfrak{J})[T_1,\ldots,T_n]\to\gr_\bullet(A)$ es biyectivo. Se concluye
que el homomorfismo canónico ($0_{\mathrm{III}}$, 10.1.1.2)
\[
  \gr_0(M)\otimes_{A/\mathfrak{J}}\gr_\bullet(A)
  \longrightarrow\gr_\bullet(M)
\]
es biyectivo. La conclusión se deduce entonces de ($0_{\mathrm{III}}$,
10.2.2), pues los $f_i$ están en el ideal máximo de $A$ y el homomorfismo
$\varphi$ es local.
\end{proof}

\subsection{Sucesiones \texorpdfstring{$\mathscr{F}$}{F}-regulares}
\label{subsection:0.15.2-es}

\begin{env}[15.2.1]
\label{0.15.2.1-es}
Sean $X$ un espacio anillado, $\mathscr{F}$ un $\mathcal{O}_X$-Módulo,
$f_i$ ($1\leq i\leq n$) una sucesión de secciones de $\mathcal{O}_X$ sobre
$X$, y $\mathscr{J}$ el Ideal de $\mathcal{O}_X$ generado por los $f_i$
(imagen de $\mathcal{O}_X^n$ por el homomorfismo
$\mathcal{O}_X^n\to\mathcal{O}_X$ definido canónicamente por los $f_i$
($0_{\mathrm I}$, 5.1.1)). Los sub-$\mathcal{O}_X$-Módulos
$\mathscr{J}^k\mathscr{F}$ definen una filtración sobre $\mathscr{F}$, y
se pone también
$\gr_k(\mathscr{F})=\mathscr{J}^k\mathscr{F}/\mathscr{J}^{k+1}\mathscr{F}$~;
si $S_k$ es el conjunto de los polinomios homogéneos de grado total $k$ en
el anillo $\mathbb{Z}[T_1,\ldots,T_n]$, $S_k$ puede considerarse como un
haz simple sobre $X$, y se define, como en (15.1.1), un homomorfismo
canónico
\[
\label{0.15.2.1.1-es}
  \varphi_k:\gr_0(\mathscr{F})\otimes_{\mathbb Z}S_k
  \longrightarrow\gr_k(\mathscr{F}).
\tag{15.2.1.1}
\]
\oldpage[0\textsubscript{IV-1}]{21}
de la manera siguiente: para todo abierto $U$ de $X$ (o de una base de la
topología de $X$), se considera el homomorfismo (15.1.1.1)
\[
  \varphi_{k,U}:\gr_0(\Gamma(U,\mathscr{F}))\otimes_{\mathbb Z}S_k
  \longrightarrow\gr_k(\Gamma(U,\mathscr{F}))
\]
definido por las $n$ secciones $f_i|U\in\Gamma(U,\mathcal{O}_X)$, estando el
$\Gamma(U,\mathcal{O}_X)$-módulo $\Gamma(U,\mathscr{F})$ filtrado por los
submódulos $(\Gamma(U,\mathscr{J}))^k\Gamma(U,\mathscr{F})$. Como el
$\mathcal{O}_X$-Módulo $\gr_k(\mathscr{F})$ está asociado al prehaz
$U\mapsto\gr_k(\Gamma(U,\mathscr{F}))$ para todo $k$, los $\varphi_{k,U}$
definen un homomorfismo de $\mathcal{O}_X$-Módulos (15.2.1.1). En virtud
de las propiedades de exactitud de los límites inductivos de módulos y de su
conmutación con los productos tensoriales, de lo anterior resulta que, para
todo $x\in X$, el homomorfismo $(\varphi_k)_x$ de los tallos en el punto $x$
de los dos miembros de (15.2.1.1) se identifica con el homomorfismo de (15.1.1)
\[
\label{0.15.2.1.2-es}
  \gr_0(\mathscr{F}_x)\otimes_{\mathbb Z}S_k
  \longrightarrow\gr_k(\mathscr{F}_x)
\tag{15.2.1.2}
\]
definido por la sucesión $((f_i)_x)_{1\leq i\leq n}$ de elementos de
$\mathcal{O}_x$.
\end{env}

\begin{env}[15.2.2]
\label{0.15.2.2-es}
Con la notación de (15.2.1), diremos que la sucesión
$(f_i)_{1\leq i\leq n}$ es $\mathscr{F}$-regular si, denotando por
$f_j\mathscr{F}$ el sub-$\mathcal{O}_X$-Módulo de $\mathscr{F}$ imagen del
endomorfismo de $\mathscr{F}$ definido por $f_j$, el endomorfismo de
$\mathscr{F}/(\sum_{j=1}^{i-1}f_j\mathscr{F})$ definido por $f_i$ es inyectivo
para $1\leq i\leq n$; diremos que la sucesión $(f_i)$ es
$\mathscr{F}$-cuasirregular si los homomorfismos (15.2.1.1) son inyectivos.
De estas definiciones y de las observaciones de (15.2.1) resulta que, para que
la sucesión $(f_i)$ sea $\mathscr{F}$-regular (resp.
$\mathscr{F}$-cuasirregular), es necesario y suficiente que, para todo
$x\in X$, la sucesión $((f_i)_x)_{1\leq i\leq n}$ sea
$\mathscr{F}_x$-regular (resp. $\mathscr{F}_x$-cuasirregular). Toda sucesión
$\mathscr{F}$-regular es, pues, $\mathscr{F}$-cuasirregular (15.1.9); el
recíproco es cierto si $\mathscr{F}$ es un $\mathcal{O}_X$-Módulo de tipo
finito, si los $\mathcal{O}_x$ son anillos noetherianos para todo $x\in X$ y
si, además, los $(f_i)_x$ pertenecen al radical de $\mathcal{O}_x$ para todo
$x\in X$ (15.1.11).

Si $X=\operatorname{Spec}(A)$ es un esquema afín, y
$\mathscr{F}=\widetilde{M}$, resulta de (I, 1.3) que, para que la sucesión
$(f_i)$ sea $\mathscr{F}$-regular (resp. $\mathscr{F}$-cuasirregular), es
necesario y suficiente que sea $M$-regular (resp. $M$-cuasirregular).

Cuando $\mathscr{F}=\mathcal{O}_X$, se omitirá la mención de $\mathscr{F}$ en
las definiciones precedentes.
\end{env}

\begin{remark}[15.2.3]
\label{0.15.2.3-es}
Cuando $X=\operatorname{Spec}(A)$ y
$\mathfrak{J}=\sum_{i=1}^n f_iA$, la sucesión $(f_j)$ es cuasirregular para
todo $A$-módulo $M$ en todo abierto $U$ que no interseque
$V(\mathfrak{J})$, pues en todo punto $x\in U$ se tiene
$\mathfrak{J}_x=\mathcal{O}_x$ y los dos miembros de (15.2.1.2) son nulos. La
noción de sucesión $\widetilde{M}$-regular solo tiene interés en un entorno de
los puntos de $V(\mathfrak{J})$.
\end{remark}

\begin{proposition}[15.2.4]
\label{0.15.2.4-es}
Sean $X$ un espacio anillado, $\mathscr{F}$ un $\mathcal{O}_X$-Módulo
coherente, $x$ un punto de $X$, y $(f_i)_{1\leq i\leq n}$ una sucesión de
elementos de $\Gamma(X,\mathcal{O}_X)$. Si la sucesión
$((f_i)_x)_{1\leq i\leq n}$ es $\mathscr{F}_x$-regular, existe un entorno
abierto $U$ de $x$ tal que la sucesión $(f_i|U)_{1\leq i\leq n}$ sea
$(\mathscr{F}|U)$-regular.
\end{proposition}

\begin{proof}
En efecto, para $1\leq i\leq n$, el $\mathcal{O}_X$-Módulo
$\mathscr{F}/(\sum_{j=1}^{i-1}f_j\mathscr{F})$ es coherente
($0_{\mathrm I}$, 5.3.4), por lo que también lo es el núcleo del endomorfismo
de este $\mathcal{O}_X$-Módulo definido por $f_i$
($0_{\mathrm I}$, 5.3.4); como el tallo de este núcleo en el punto $x$ es
nulo, lo mismo ocurre con su restricción a un entorno de $x$
($0_{\mathrm I}$, 5.2.2).
\end{proof}

\begin{proposition}[15.2.5]
\label{0.15.2.5-es}
\oldpage[0\textsubscript{IV-1}]{22}
Sea $u=(\psi,\theta):X\to Y$ un morfismo plano de espacios anillados
($0_{\mathrm I}$, 6.7.1); sean $\mathscr{F}$ un
$\mathcal{O}_Y$-Módulo, $\mathscr{G}=u^*(\mathscr{F})$,
$(f_i)_{1\leq i\leq n}$ una sucesión de elementos de
$\Gamma(Y,\mathcal{O}_Y)$, y $(g_i)_{1\leq i\leq n}$ la sucesión de sus
imágenes
$g_i=\theta(f_i)\in\Gamma(Y,u_*(\mathcal{O}_X))=\Gamma(X,\mathcal{O}_X)$. Si
la sucesión $(f_i)$ es $\mathscr{F}$-regular, la sucesión $(g_i)$ es
$\mathscr{G}$-regular. El recíproco es cierto si $u$ es un morfismo fielmente
plano ($0_{\mathrm I}$, 6.7.8).
\end{proposition}

\begin{proof}
Sean $x$ un punto de $X$ e $y=u(x)=\psi(x)$; se tiene
\[
  \mathscr{G}_x=\psi_x^*(\mathscr{F}_y)
  \otimes_{\psi_x^*(\mathcal{O}_y)}\mathcal{O}_x,
  \qquad\text{y}\qquad
  (g_i)_x=\theta_x^\#(\psi_x^*((f_i)_y))~;
\]
como, por hipótesis,
$\theta_x^\#:\psi_x^*(\mathcal{O}_y)\to\mathcal{O}_x$ hace de
$\mathcal{O}_x$ un $\psi_x^*(\mathcal{O}_y)$-módulo plano, y como el funtor
$\psi^*$ es exacto, la primera afirmación resulta de la afirmación
correspondiente de (15.1.14); lo mismo ocurre con la segunda, teniendo en cuenta
que la hipótesis de que $u$ es fielmente plano significa que $\psi$ es
sobreyectivo y que $\theta_x^\#$ hace de $\mathcal{O}_x$ un
$\psi_x^*(\mathcal{O}_y)$-módulo fielmente plano.
\end{proof}
